% Apéndice: pentágono regular: cuentas y script
\section{Pentágono regular: cuentas y script}\label{cap:18}\label{ap:pentagono}

Ternas homogéneas (cuatro decimales):
\[
A=\begin{pmatrix}0\\1\\1\end{pmatrix},\
B=\begin{pmatrix}-0.9511\\0.3090\\1\end{pmatrix},\
C=\begin{pmatrix}-0.5878\\-0.8090\\1\end{pmatrix},\
D=\begin{pmatrix}0.5878\\-0.8090\\1\end{pmatrix},\
E=\begin{pmatrix}0.9511\\0.3090\\1\end{pmatrix}.
\]

\subsection{Verificación de los tres casos}
\begin{center}\small
\begin{tabularx}{\textwidth}{>{\bfseries}l X X X}
\toprule
& (a) \(ABCD\to BCDE\) & (b) \(ABCD\to BACD\) & (c) \(ABCD\to CBAD\)\\
\midrule
\(HA\) & \(B\), \(W=1\) & \(B\), \(W=9.4721\) & \(C\), \(W=1.3820\)\\
\(HB\) & \(C\), \(W=1\) & \(A\), \(W=3.6180\) & \(B\), \(W=0.8541\)\\
\(HC\) & \(D\), \(W=1\) & \(C\), \(W=-5.8541\) & \(A\), \(W=0.5279\)\\
\(HD\) & \(E\), \(W=1\) & \(D\), \(W=-5.8541\) & \(D\), \(W=0.8541\)\\
\(HE\) & \(A\), \(W=1\) & \((-3.0777,1)\), \(W=3.6180\) & \((0,-1.2361)\), \(W=1.3820\)\\
recta al infinito & ideal & \(y=-0.1180\) & \(0.2775x+0.3820y+1=0\)\\
autovalores & \(1,\ e^{\pm i72^\circ}\) & \(5.8541,\ -5.8541\ (\times2)\) &
\(-0.8541,\ 0.8541\ (\times2)\)\\
puntos fijos & \(O=(0,0)\) & eje \(CD\); centro \((-0.5878,0.5729)\) &
eje \(BD\); centro \((-1.5388,-3.7361)\)\\
tipo & rotación (afín) & homología armónica, imagen partida & homología armónica\\
\bottomrule
\end{tabularx}
\end{center}

\subsection{Autovector por producto cruz, caso (c)}
Polinomio característico de \(H_c\): \(\lambda^3-0.8541\lambda^2-0.7295\lambda+0.6231\),
raíces \(0.8541\) (doble) y \(-0.8541\). Para \(\lambda=-0.8541\),
\[
H_c+0.8541\,I=\begin{pmatrix}1.2812&-0.5878&-0.2245\\-1.0368&0.2812&-0.5451\\0.2775&0.3820&1.8541\end{pmatrix},
\]
\[
\text{fila}_1\times\text{fila}_2=
\begin{pmatrix}(-0.5878)(-0.5451)-(-0.2245)(0.2812)\\(-0.2245)(-1.0368)-(1.2812)(-0.5451)\\(1.2812)(0.2812)-(-0.5878)(-1.0368)\end{pmatrix}
=\begin{pmatrix}0.3835\\0.9311\\-0.2492\end{pmatrix}\sim\begin{pmatrix}-1.5388\\-3.7361\\1\end{pmatrix}.
\]
Comprobación:
\[
  H_c\begin{pmatrix}-1.5388\\-3.7361\\1\end{pmatrix}=\begin{pmatrix}1.3143\\3.1909\\-0.8541\end{pmatrix}
  =-0.8541\begin{pmatrix}-1.5388\\-3.7361\\1\end{pmatrix}.
\]
Para \(\lambda=0.8541\) el producto cruz de cualquier par de filas de \(H_c-0.8541I\) da
\((0,0,0)\): raíz doble, toda la recta \(BD\) es de puntos fijos.

\subsection{Imagen partida, caso (b)}
\begin{center}
\begin{tikzpicture}[scale=1.15]
  \clip (-3.6,-2.2) rectangle (2.4,2.1);
  \draw[linegray!40,very thin] (-3.5,-2) grid[step=.5] (2,2);
  \draw[axisline,-{Latex[length=2.2mm]}] (-3.5,0)--(2.2,0) node[right] {$x$};
  \draw[axisline,-{Latex[length=2.2mm]}] (0,-2.1)--(0,2) node[above] {$y$};
  \draw[teal,line width=1pt,fill=teallight!45] (0,1)--(-0.9511,0.3090)--(-0.5878,-0.8090)--(0.5878,-0.8090)--(0.9511,0.3090)--cycle;
  \foreach \p/\n/\pos in {(0,1)/A/{above right},(-0.9511,0.3090)/B/{above left},(-0.5878,-0.8090)/C/{below left},(0.5878,-0.8090)/D/{below right},(0.9511,0.3090)/E/right}
    \fill[teal] \p circle (1.2pt) node[\pos,font=\small,text=tealdark] {$\n$};
  \draw[gold!80!black,dotted,line width=.9pt] (-3.6,-0.118)--(2.4,-0.118) node[pos=.08,above,font=\scriptsize,text=gold!60!black] {$y=-0.118$ (se va al infinito)};
  % imagen: segmentos AB, CD, E'B
  \draw[rust,line width=1pt] (0,1)--(-0.9511,0.3090);
  \draw[rust,line width=1pt] (-0.5878,-0.8090)--(0.5878,-0.8090);
  \draw[rust,line width=1pt] (-3.0777,1)--(-0.9511,0.3090);
  % imagen de BC: recta AC exterior al segmento AC
  \draw[rust,line width=1pt,-{Latex[length=1.8mm]}] (0,1)--(0.35,2.077);
  \draw[rust,line width=1pt,-{Latex[length=1.8mm]}] (-0.5878,-0.8090)--(-1.0,-2.077);
  % imagen de DE: recta D E' exterior al segmento
  \draw[rust,line width=1pt,-{Latex[length=1.8mm]}] (0.5878,-0.8090)--(2.4,-1.703);
  \draw[rust,line width=1pt,-{Latex[length=1.8mm]}] (-3.0777,1)--(-3.6,1.258);
  \fill[rust] (-3.0777,1) circle (1.4pt) node[above,font=\small,text=rust,labelback] {$E'_b=(-3.08,\,1)$};
  \draw[dashed,tealdark,line width=.8pt] (-3.6,-0.809)--(2.4,-0.809) node[pos=.12,below,font=\scriptsize,text=tealdark] {eje $CD$ (puntos fijos)};
  \fill[ink] (-0.5878,0.5729) circle (1.4pt) node[left,font=\scriptsize,labelback] {$O_b$};
\end{tikzpicture}
\end{center}
Los lados \(AB\), \(CD\) y \(EA\) van a los segmentos \(BA\), \(CD\) y \(E'_bB\). Los lados
\(BC\) y \(DE\) cruzan la recta \(y=-0.118\) y sus imágenes son las semirrectas exteriores de
las rectas \(AC\) y \(DE'_b\).

\subsection{Script}
Reproduce todo lo anterior: patrón, homografía, imágenes, autovalores, autovector por
producto cruz de filas y verificación.
\begin{verbatim}
import numpy as np
np.set_printoptions(precision=4, suppress=True)
ang = np.deg2rad(90 + 72*np.arange(5))
V = np.c_[np.cos(ang), np.sin(ang), np.ones(5)]      # A..E como filas
A, B, C, D, E = V

def patron(P):                       # P: 4 ternas como columnas
    M = P[:, :3]; lam = np.linalg.solve(M, P[:, 3]); return M * lam
def homografia(src, dst):
    K1 = patron(np.column_stack(src)); K2 = patron(np.column_stack(dst))
    H = K2 @ np.linalg.inv(K1)
    return H / H[2, 2]

casos = {"a": (B, C, D, E), "b": (B, A, C, D), "c": (C, B, A, D)}
for k, dst in casos.items():
    H = homografia((A, B, C, D), dst); print("caso", k); print(H)
    w = H @ E; print(" E ->", w[:2] / w[2], " W =", round(w[2], 4))
    lams = np.linalg.eigvals(H); print(" autovalores", lams)
    for lam in lams[np.isreal(lams)].real:
        M = H - lam * np.eye(3)
        v = np.cross(M[0], M[1])        # producto cruz de dos filas
        if np.linalg.norm(v) < 1e-9:
            v = np.cross(M[0], M[2])
        if np.linalg.norm(v) < 1e-9:
            print(f"  lambda={lam:.4f}: raiz doble, recta fija")
            continue
        if abs(v[2]) > 1e-9:
            v = v / v[2]                # punto propio (x, y, 1)
        print(f"  lambda={lam:.4f}: punto fijo", v)
\end{verbatim}

